\documentclass[12pt]{article}
\usepackage[utf8]{inputenc}
\usepackage[T1]{fontenc}
\usepackage{hyperref}
\title{Uzbek (macrolanguage) Language Resources for K-12 Curriculum Development}
\author{Compiled by Claude Sonnet 5 (claude-sonnet-5) for Andr\'as Kornai}
\date{}
\begin{document}
\maketitle

\section{Introduction}
The code uzb is the ISO 639-3 macrolanguage code for Uzbek. It includes Northern Uzbek (uzn, spoken mainly in Uzbekistan) and Southern Uzbek (uzs, spoken mainly in Afghanistan and Pakistan). Search results gave unstable speaker counts, and I do not rely on any. Most resources found concern Northern Uzbek; see the separate uzn dossier for school materials. Nothing found was specific to uzb as distinct from its member languages.

\section{Primary Educational Resources}
Resources found are mostly for foreign learners of Uzbek. Government school textbooks are covered in the uzn dossier.
\begin{itemize}
\item Georgetown University Press, Uzbek textbook front matter (Azimova, elementary, 2010): \url{https://press.georgetown.edu/files/materials/9781589017061_Frontmatter.pdf}
\item Cambridge Language Centre catalogue entry, Uzbek: an elementary textbook / O'zbek tili: birinchi bosqich: \url{https://www.langcen.cam.ac.uk/resources/langu/data/uz_108.html}
\end{itemize}

\section{Linguistic and Reference Materials}
\begin{itemize}
\item Wikipedia, Uzbek language (describes the macrolanguage and its two main variants): \url{https://en.wikipedia.org/wiki/Uzbek_language}
\item University of Chicago CEERES, Uzbek Language Resources (links to online dictionaries): \url{https://ceeres.uchicago.edu/uzbek-language-resources}
\end{itemize}

\section{Religious and Community Resources}
None found in this search.

\section{Digital and Audio Resources}
\begin{itemize}
\item Omniglot Uzbek book listing: \url{https://www.omniglot.com/books/language/uzbek.htm}
\end{itemize}
No dedicated Uzbek corpus was found.

\section{Curriculum Development Context}
Uzbek is written in both Latin and Cyrillic scripts. The Georgetown intermediate textbook uses Cyrillic, and government reports and media are described as still using it. A curriculum under the code uzb should specify which variety it targets, since uzn and uzs differ in script and setting.
\end{document}